\SourceLeaf{212}
\Needspace{12\baselineskip}
\begin{center}
\textit{Formen, die unter anderen enthalten sind, ohne ihnen äquivalent zu sein.}
\end{center}

\Needspace{8\baselineskip}
\Article{213}
Oben in den Artikeln~157, 158 haben wir gezeigt: Enthält eine Form $f$ mit der Determinante $D$ eine Form $F$ mit der Determinante $E$ und geht sie durch die Substitution $\alpha,\beta,\gamma,\delta$ in diese über, so ist $E=(\alpha\delta-\beta\gamma)^2D$. Ist $\alpha\delta-\beta\gamma=\pm1$, so enthält die Form $f$ nicht nur $F$, sondern ist ihr auch äquivalent. Wenn also $f$ die Form $F$ enthält, ihr aber nicht äquivalent ist, so ist der Quotient $\frac{E}{D}$ eine ganze Zahl größer als 1. Wir müssen daher jetzt \textit{entscheiden, ob eine gegebene Form $f$ mit der Determinante $D$ eine gegebene Form $F$ mit der Determinante $Dee$ enthält}, wobei $e$ als positive Zahl größer als 1 vorausgesetzt wird. Dazu zeigen wir, wie sich endlich viele unter $f$ enthaltene Formen so angeben lassen, dass $F$, falls sie unter $f$ enthalten ist, notwendig einer von ihnen äquivalent sein muss.

I. Alle positiven Teiler von $e$ (einschließlich 1 und $e$) seien $m,m',m''$ usw., und es sei $e=mn=m'n'=m''n''$ usw. Der Kürze wegen bezeichnen wir die Form, in die $f$ durch die eigentliche Substitution $m.\,0,0,n$ übergeht, mit $(m;0)$, die Form, in die $f$ durch die eigentliche Substitution $m,1,0,n$ übergeht, mit $(m;1)$ usw. Allgemein bezeichnen wir die Form, in die $f$ durch die eigentliche Substitution $m,k,0,n$ übergeht, mit $(m;k)$. Ebenso gehe $f$ durch die eigentliche Substitution $m',0,0,n'$ in $(m';0)$, durch $m',1,0,n'$ in $(m';1)$ usw., durch $m'',0,0,n''$ in $(m'';0)$ usw. über. Alle diese Formen sind unter $f$ eigentlich enthalten, und jede hat die Determinante $=Dee$. Die Gesamtheit aller Formen $(m;0)$, $(m;1)$, $(m;2)\ldotp\ldotp\ldotp\ldotp(m;m-1)$, $(m';0)$, $(m';1)\ldots(m';m'-1)$; $(m'';0)$ usw., deren Anzahl $m+m'+m''+\text{usw.}$ ist und die offenbar alle verschieden sind, bezeichnen wir mit $\Omega$.

Ist \textit{z.B.} die Form $f$ gleich $(2,5,7)$ und gilt $e=5$, so enthält $\Omega$ die folgenden sechs Formen: $(1;0)$, $(5;0)$, $(5;1)$, $(5;2)$, $(5;3)$, $(5;4)$. Ausgeschrieben sind dies $(2,25,175)$, $(50,25,7)$, $(50,35,19)$, $(50,45,35)$, $(50,55,55)$, $(50,65,79)$.

II. Ich behaupte nun: Ist eine Form $F$ mit der Determinante $Dee$ unter $f$ eigentlich enthalten, so muss sie einer der Formen $\Omega$ eigentlich äquivalent sein. Nehmen wir an, die Form $f$ gehe in $F$ über, und zwar durch die eigentliche Substitution $\alpha,\beta,\gamma,\delta$. Dann ist $\alpha\delta-\beta\gamma=e$. Der größte gemeinsame Teiler der Zahlen $\gamma,\delta$ (die nicht beide 0 sein können), positiv genommen, sei $=n$, und es sei $\frac{e}{n}=m$, was offenbar eine ganze
\SourceLeaf{213}
Zahl ist. Wir wählen $g,h$ so, dass $\gamma g+\delta h=n$ ist, und schließlich sei $k$ der kleinste positive Rest von $\alpha g+\beta h$ modulo $m$. Dann ist $(m;k)$, die offenbar zu den Formen $\Omega$ gehört, der Form $F$ eigentlich äquivalent und geht durch die eigentliche Substitution
\[
\frac{\gamma}{n}\cdot\frac{\alpha g+\beta h-k}{m}+h,\quad
\frac{\delta}{n}\cdot\frac{\alpha g+\beta h-k}{m}-g,\quad
\frac{\gamma}{n},\quad\frac{\delta}{n}
\]
in sie über. Denn \textit{erstens} sind diese vier Zahlen offenbar ganze Zahlen; \textit{zweitens} lässt sich leicht nachprüfen, dass die Substitution eigentlich ist; \textit{drittens} ist die Form, in die $(m;k)$ durch diese Substitution übergeht, offenbar dieselbe wie die, in die $f$\HistoricalNote{*)}{Sie geht nämlich durch die Substitution $m,k,0,n$ in $(m;k)$ über. Siehe Artikel~159.} durch die Substitution
\[
m\left(\frac{\gamma}{n}\cdot\frac{\alpha g+\beta h-k}{m}+h\right)+\frac{k\gamma}{n},\quad
m\left(\frac{\delta}{n}\cdot\frac{\alpha g+\beta h-k}{m}-g\right)+\frac{k\delta}{n},\quad\gamma,\quad\delta
\]
übergeht, oder wegen $mn=e=\alpha\delta-\beta\gamma$ und daher $\beta\gamma+mn=\alpha\delta$, $\alpha\delta-mn=\beta\gamma$ durch
\[
\frac{1}{n}(\alpha\gamma g+\alpha\delta h),\quad
\frac{1}{n}(\beta\gamma g+\beta\delta h),\quad\gamma,\quad\delta
\]
oder schließlich wegen $\gamma g+\delta h=n$ durch $\alpha,\beta,\gamma,\delta$, \textit{d.h.} nach Voraussetzung in $F$. Also sind $(m;k)$ und $F$ eigentlich äquivalent. \textit{Damit ist die Behauptung bewiesen.}

So lässt sich immer entscheiden, ob eine gegebene Form $f$ mit der Determinante $D$ eine Form $F$ mit der Determinante $Dee$ eigentlich enthält. Wollen wir entscheiden, ob $f$ die Form $F$ uneigentlich enthält, brauchen wir nur zu untersuchen, ob die zu $F$ entgegengesetzte Form unter $f$ eigentlich enthalten ist (Artikel~159).

\Needspace{8\baselineskip}
\Article{214}
\textsc{Aufgabe.} \textit{Für zwei gegebene Formen, $f$ mit der Determinante $D$ und $F$ mit der Determinante $Dee$, von denen die erste die zweite eigentlich enthält, alle eigentlichen Transformationen von $f$ in $F$ angeben.}

\textit{Lösung.} Mit $\Omega$ bezeichnen wir dieselbe Gesamtheit von Formen wie im vorhergehenden Artikel. Daraus wählen wir alle Formen aus, denen $F$ eigentlich äquivalent ist; sie seien $\Phi,\Phi',\Phi''$ usw. Jede dieser Formen liefert auf die folgende Weise eigentliche Transformationen von $f$ in $F$, wobei verschiedene Formen verschiedene Transformationen liefern (\textit{d.h.} jede Form liefert andere), alle zusammen aber sämtliche (\textit{d.h.} es gibt keine eigentliche Transformation von $f$ in $F$, die nicht von einer der Formen $\Phi,\Phi'$ usw. geliefert wird). Da die Methode für alle Formen $\Phi,\Phi'$ usw. dieselbe ist, sprechen wir nur von einer.

\SourceLeaf{214}
Es sei $\Phi$ gleich $(M;K)$ und $e=MN$, sodass $f$ in $\Phi$ übergeht, und zwar durch die eigentliche Substitution $M,K,0,N$. Ferner bezeichnen wir alle eigentlichen Transformationen von $\Phi$ in $F$ mit $\mathfrak{a},\mathfrak{b},\mathfrak{c},\mathfrak{d}$. Dann geht offenbar $f$ in $\Phi$ über, und zwar durch die eigentliche Substitution $M\mathfrak{a}+K\mathfrak{c},M\mathfrak{b}+K\mathfrak{d},N\mathfrak{c},N\mathfrak{d}$. Auf diese Weise liefert jede eigentliche Transformation von $\Phi$ in $F$ eine eigentliche Transformation von $f$ in $F$. --- Die übrigen Formen $\Phi',\Phi''$ usw. behandeln wir genauso: Jede ihrer eigentlichen Transformationen in $F$ liefert eine eigentliche Transformation von $f$ in $F$.

Um zu sehen, dass diese Lösung in jeder Hinsicht vollständig ist, müssen wir zeigen:

I. \textit{Auf diese Weise erhalten wir alle möglichen eigentlichen Transformationen von $f$ in $F$.} Eine beliebige eigentliche Transformation von $f$ in $F$ sei $\alpha,\beta,\gamma,\delta$, und wie in II des vorhergehenden Artikels sei $n$ der größte gemeinsame Teiler von $\gamma,\delta$. Die Zahlen $m,g,h,k$ bestimmen wir wie dort. Dann ist $(m;k)$ unter den Formen $\Phi,\Phi'$ usw. enthalten, und
\[
\frac{\gamma}{n}\cdot\frac{\alpha g+\beta h-k}{m}+h,\quad
\frac{\delta}{n}\cdot\frac{\alpha g+\beta h-k}{m}-g,\quad
\frac{\gamma}{n},\quad\frac{\delta}{n}
\]
ist eine ihrer eigentlichen Transformationen in $F$. Daraus erhalten wir nach der eben angegebenen Regel die Transformation $\alpha,\beta,\gamma,\delta$. Dies alles wurde im vorhergehenden Artikel bewiesen.

\Needspace{26\baselineskip}
II. \textit{Alle auf diese Weise entstehenden Transformationen sind verschieden; keine wird zweimal erhalten.} Wir sehen leicht, dass verschiedene Transformationen derselben Form $\Phi$ oder $\Phi'$ usw. in $F$ nicht dieselbe Transformation von $f$ in $F$ erzeugen können. Dass aber auch verschiedene Formen, \textit{z.B.} $\Phi$ und $\Phi'$, nicht dieselbe Transformation liefern können, beweisen wir so. Nehmen wir an, die eigentliche Transformation $\alpha,\beta,\gamma,\delta$ von $f$ in $F$ werde \textit{sowohl} aus der eigentlichen Transformation $\mathfrak{a},\mathfrak{b},\mathfrak{c},\mathfrak{d}$ von $\Phi$ in $F$ \textit{als auch} aus der eigentlichen Transformation $\mathfrak{a}',\mathfrak{b}',\mathfrak{c}',\mathfrak{d}'$ von $\Phi'$ in $F$ erhalten. Es sei $\Phi=(M;K)$, $\Phi'=(M';K')$, $e=MN=M'N'$. Dann haben wir die Gleichungen
\begin{align*}
\alpha&=M\mathfrak{a}+K\mathfrak{c}=M'\mathfrak{a}'+K'\mathfrak{c}' \tag*{. . . . . . . [1]}\\
\beta&=M\mathfrak{b}+K\mathfrak{d}=M'\mathfrak{b}'+K'\mathfrak{d}' \tag*{. . . . . . . [2]}\\
\gamma&=N\mathfrak{c}=N'\mathfrak{c}' \tag*{. . . . . . . . . . . . . [3]}
\\
\SourceLeaf{215}\delta&=N\mathfrak{d}=N'\mathfrak{d}' \tag*{. . . . . . . . . . . . . [4]}\\
\mathfrak{a}\mathfrak{d}-\mathfrak{b}\mathfrak{c}&=\mathfrak{a}'\mathfrak{d}'-\mathfrak{b}'\mathfrak{c}'=1 \tag*{. . . . . . . . . [5]}
\end{align*}
Aus $\mathfrak{a}[4]-\mathfrak{b}[3]$ folgt mithilfe von Gleichung [5] $N=N'(\mathfrak{a}\mathfrak{d}'-\mathfrak{b}\mathfrak{c}')$, also teilt $N'$ die Zahl $N$. Ebenso ergibt $\mathfrak{a}'[4]-\mathfrak{b}'[3]$ die Gleichung $N(\mathfrak{a}'\mathfrak{d}-\mathfrak{b}'\mathfrak{c})=N'$, also teilt $N$ die Zahl $N'$. Da sowohl $N$ als auch $N'$ als positiv vorausgesetzt werden, muss $N=N'$ und $M=M'$ sein. Aus [3] und [4] folgt dann $\mathfrak{c}=\mathfrak{c}'$, $\mathfrak{d}=\mathfrak{d}'$. Ferner folgt aus $\mathfrak{a}[2]-\mathfrak{b}[1]$
\[
K=M'(\mathfrak{a}\mathfrak{b}'-\mathfrak{b}\mathfrak{a}')+K'(\mathfrak{a}\mathfrak{d}'-\mathfrak{b}\mathfrak{c}')=M(\mathfrak{a}\mathfrak{b}'-\mathfrak{b}\mathfrak{a}')+K'
\]
also $K\equiv K'\;(\mathrm{mod.}\,M)$. Das ist nur möglich, wenn $K=K'$ ist, denn sowohl $K$ als auch $K'$ liegen zwischen 0 und $M-1$. Daher sind die Formen $\Phi,\Phi'$ nicht verschieden, entgegen der Voraussetzung.

Ist $D$ negativ oder positiv quadratisch, so lassen sich mit dieser Methode offenbar alle eigentlichen Transformationen von $f$ in $F$ tatsächlich finden. Ist dagegen $D$ positiv und nichtquadratisch, so können wir allgemeine Formeln angeben, die alle eigentlichen Transformationen enthalten, deren Anzahl unendlich ist.

Ist schließlich die Form $F$ unter der Form $f$ uneigentlich enthalten, so lassen sich nach der angegebenen Methode alle uneigentlichen Transformationen der ersten in die zweite leicht darstellen. Sollen nämlich $\alpha,\beta,\gamma,\delta$ alle eigentlichen Transformationen von $f$ in die zu $F$ entgegengesetzte Form durchlaufen, so werden alle uneigentlichen Transformationen von $f$ in $F$ durch $\alpha,-\beta,\gamma,-\delta$ dargestellt.

\textit{Beispiel.} Gesucht sind alle Transformationen von $(2,5,7)$ in $(275,0,-1)$, die sowohl eigentlich als auch uneigentlich unter der ersten enthalten ist. Die Gesamtheit $\Omega$ für diesen Fall haben wir schon im vorhergehenden Artikel angegeben. Bei der Untersuchung finden wir, dass sowohl $(5;1)$ als auch $(5;4)$ der Form $(275,0,-1)$ eigentlich äquivalent sind. Nach unserer oben entwickelten Theorie sind alle eigentlichen Transformationen von $(5;1)$, \textit{d.h.} $(50,35,19)$, in $(275,0,-1)$ in der allgemeinen Formel
\[
16t-275u,\quad -t+16u,\quad -15t+275u,\quad t-15u
\]
enthalten, wobei $t,u$ alle ganzen Zahlen durchlaufen, die $tt-275uu=1$
\SourceLeaf{216}
erfüllen. Daher sind alle daraus entstehenden eigentlichen Transformationen von $(2,5,7)$ in $(275,0,-1)$ in der allgemeinen Formel
\[
65t-1100u,\quad -4t+65u,\quad -15t+275u,\quad t-15u
\]
enthalten. Ebenso sind alle eigentlichen Transformationen von $(5,4)$, \textit{d.h.} $(50,65,79)$, in $(275,0,-1)$ in der allgemeinen Formel
\[
14t+275u,\quad t+14u,\quad -15t-275u,\quad -t-15u
\]
enthalten und somit alle daraus entstehenden eigentlichen Transformationen von $(2,5,7)$ in $(275,0,-1)$ in
\[
10t+275u,\quad t+10u,\quad -15t-275u,\quad -t-15u
\]
enthalten. Diese beiden Formeln umfassen also alle gesuchten eigentlichen Transformationen\HistoricalNote{*)}{Eleganter werden alle eigentlichen Transformationen durch die Formel
\[
10t+55u,\quad t+2u,\quad -15t-55u,\quad -t-3u
\]
dargestellt, wobei $t,u$ alle ganzen Zahlen durchlaufen, die $tt-11uu=1$ erfüllen.}. --- Auf dieselbe Weise finden wir, dass alle uneigentlichen Transformationen von $(2,5,7)$ in $(275,0,-1)$ in den folgenden beiden Formeln enthalten sind:
\[
\text{(I)}\;\ldots\quad 65t-1100u,\quad 4t-65u,\quad -15t+275u,\quad -t+15u
\]
und
\[
\text{(II)}\;\ldots\quad 10t+275u,\quad -t-10u,\quad -15t-275u,\quad t+15u
\]

